% #############################################################################
% This is Appendix A
% !TEX root = ../main.tex
% #############################################################################
\chapter{Additional Data Preparation Details}
\label{chap:appendix-data-details}
\label{chap:appendix-opensubs-filtering}

This appendix provides additional details on the preparation of the resources
used in the experiments, focusing on OpenSubtitles and GPT-Wikipedia, whose
construction involved substantial preprocessing and filtering. FRMT and the Golden Collection mainly require conversion to the
common experimental format and are described in
\Cref{sec:data-frmt,sec:data-golden-collection}. This appendix
also reports tokenization statistics for the resources used in the experiments.

\section{Tokenization Statistics}
\label{sec:appendix-tokenization-statistics}

\looseness=-1
\Cref{tab:appendix-tokenization-statistics} reports corpus token counts and the
number of tokenizer vocabulary entries observed with the
\texttt{google/t5gemma-2-4b-4b} tokenizer used by the 4B-parameter T5Gemma 2
model.
These statistics are computed over the complete processed resource splits shown
in the table, independently of how many examples are subsequently presented to
a model during training.
\textit{Total tokens} counts every tokenizer-token occurrence across all
sentences in the specified resource variant, whereas \textit{Distinct
tokenizer identifiers (IDs)} counts each vocabulary ID only once, regardless of how often it
occurs. The large difference between the quantities reported in Table A.1 is
therefore expected, since a corpus can contain millions of token occurrences
while repeatedly drawing from a much smaller set of vocabulary entries.

\begin{table}[htbp]
\centering
\thesistablefont
\setlength{\tabcolsep}{5pt}
\caption{Token and tokenizer-vocabulary statistics for the resources used in the reported experiments.}
\label{tab:appendix-tokenization-statistics}
\begin{tabularx}{\textwidth}{@{}lYccc@{}}
\toprule
Resource & Split and variety & Rows & Total tokens & \shortstack{Distinct\\tokenizer IDs} \\
\midrule
OpenSubtitles & train, pt-BR & 10,347,883 & 113,121,931 & 73,773 \\
OpenSubtitles & train, pt-PT & 10,347,883 & 116,213,290 & 74,179 \\
\midrule
GPT-Wikipedia & train, pt-BR & 4,519 & 121,589 & 8,166 \\
GPT-Wikipedia & train, pt-PT & 4,519 & 122,542 & 8,127 \\
\midrule
FRMT & train, pt-BR & 2,502 & 94,966 & 11,044 \\
FRMT & train, pt-PT & 2,502 & 98,345 & 10,886 \\
FRMT & test, pt-BR & 2,611 & 97,920 & 11,117 \\
FRMT & test, pt-PT & 2,611 & 101,039 & 10,874 \\
\midrule
Golden Collection & test, pt-BR & 500 & 14,166 & 4,234 \\
Golden Collection & test, pt-PT & 500 & 14,285 & 4,209 \\
\bottomrule
\end{tabularx}
\end{table}

\section{OpenSubtitles}
\label{sec:appendix-opensubs}

OpenSubtitles is used as a large-scale source of aligned EP/BP sentence pairs.
Its main value is coverage: it exposes the models to many lexical contexts,
short utterances, informal dialogue, and everyday constructions. The same
properties also make it noisy. Subtitle files often contain segmentation
artifacts, speaker marks, incomplete fragments, and weak sentence alignment.
Preparation therefore aims to keep as many useful aligned pairs as possible
while removing cases where the paired sentences no longer express the same
content.

\subsection{Formatting and Segmentation Repairs}

As summarized in \Cref{sec:data-opensubtitles}, formatting cleanup removes
subtitle-specific artifacts such as dialogue dashes, speaker labels, irregular
spacing, wrapping artifacts, unbalanced quotation marks, and fully capitalized
formatting. This subsection focuses on merging adjacent subtitle rows to restore
sentence boundaries.

\looseness=-1
Subtitle segmentation is then corrected when adjacent rows appear to belong to
the same sentence. A continuation cue is evaluated separately for the EP and BP
sentences after formatting normalization. It holds when the current row ends in
soft punctuation, comprising a comma, semicolon, colon, em dash, en dash, or
hyphen, or contains fewer than 40 characters, and the first alphabetic character
of the following row is lowercase. A merge is allowed when this cue holds for
either sentence.

A separate segmentation-mismatch cue applies when exactly one side of the
current pair contains an internal hard stop followed later by another hard stop.
This pattern suggests that material corresponding to more than one subtitle
sentence may have been grouped on only one side. Hard punctuation comprises
periods, question marks, exclamation marks, and ellipses. Rows are otherwise
kept separate when both current sentences end in hard punctuation and the first
alphabetic character of both following sentences is uppercase. The
segmentation-mismatch cue overrides this safeguard.

Whenever a merge is triggered, both paired sentences are concatenated so that
their row correspondence is preserved. The goal is to recover sentence
boundaries split inconsistently by subtitle timing rather than to create longer
examples indiscriminately.

\Cref{tab:opensubs-merge-examples} shows
examples of these repairs.

\subsection{SimAlign Configuration}

After local cleanup, OpenSubtitles pairs are filtered with SimAlign
\citep{jalili-sabet-etal-2020-simalign}. SimAlign provides token-level links
between the BP and EP sentences using contextual multilingual representations.
This is useful for EP/BP data, since valid correspondences are not always
surface-identical. A BP gerund may correspond to an EP infinitival
construction, and a lexical item in one sentence may correspond to a short
phrase in the other.

SimAlign requires three main choices. The first is the multilingual encoder
used to compute contextual embeddings. The second is the granularity of the
alignment, since links can be produced over subword pieces or over the
original words after combining the representations of their subword pieces.
The third is the extraction method that converts token similarities into
alignment links. This work uses the XLM-R-based encoder configuration,
word-level alignments, and the conservative intersection method, referred to
as \texttt{inter} in SimAlign.

The extraction method operates over a similarity matrix, where rows correspond to
tokens in the BP sentence, columns correspond to tokens in the EP sentence, and each
cell estimates how close the two contextual representations are. The
intersection method links two tokens only when each one is the other's
highest-scoring option in this matrix. Two alternative SimAlign methods were
also considered. The iterative maximum method begins from conservative
mutual-best links, then iteratively adds further high-scoring links to cover
tokens that remain unaligned, and can therefore produce more links than the
intersection method. Maximum-weight matching considers the similarity matrix
globally and selects a high-scoring one-to-one structure that maximizes the
overall alignment score, rather than requiring every selected pair to be
mutual local maxima. The intersection method was selected for filtering since
it provides a conservative core of bilingual evidence and avoids accepting
pairs through many weaker links.

\begin{quote}
\textbf{Example.} Consider BP tokens
\textit{[Ele] [est\'a] [correndo]} and EP tokens
\textit{[Ele] [est\'a] [a] [correr]}. The BP gerund
\textit{correndo} corresponds to the EP construction
\textit{a correr}. The following illustrative matrix shows the kind of
token-similarity evidence used by the extraction step.

\begin{center}
\thesistablefont
\begin{tabular}{lcccc}
\toprule
 & \textit{Ele} & \textit{est\'a} & \textit{a} & \textit{correr} \\
\midrule
\textit{Ele} & \cellcolor{lightgray}\textbf{0.94} & 0.21 & 0.08 & 0.12 \\
\textit{est\'a} & 0.19 & \cellcolor{lightgray}\textbf{0.91} & 0.42 & 0.28 \\
\textit{correndo} & 0.14 & 0.35 & 0.48 & \cellcolor{lightgray}\textbf{0.86} \\
\bottomrule
\end{tabular}
\end{center}

\looseness=-1
The selected mutual-best links are \textit{Ele}
to \textit{Ele} (\textbf{0.94}), \textit{est\'a}
to \textit{est\'a} (\textbf{0.91}), and \textit{correndo}
to \textit{correr} (\textbf{0.86}). Each score is the
highest-scoring option for both linked tokens. The \textit{correndo}
to \textit{a} link (\textbf{0.48}) is not selected: \textit{a}
may prefer \textit{correndo}, but \textit{correndo} prefers \textit{correr} at
\textbf{0.86}, so the pair is not mutual-best. Intersection therefore captures
the main gerund-to-infinitive correspondence \textit{correndo}
to \textit{correr} while leaving the EP support token \textit{a}
unaligned.
\end{quote}

\subsection{Coverage Score and Length Filtering}

\looseness=-1
To decide whether aligned pairs are kept in the final dataset, the alignment
links provided by SimAlign are converted into a coverage score over normalized
word tokens. Text is normalized with canonical Unicode normalization and case
folding. The scoring tokens keep alphabetic forms, Portuguese diacritics,
internal apostrophes or hyphens, and digit tokens. A token type is a distinct
normalized token form, regardless of how many times that form occurs in the
sentence. The coverage score is therefore type-based: repeated occurrences of
the same normalized form do not contribute additional distinct types.

Let \(T_{\mathrm{ep}}\) and \(T_{\mathrm{bp}}\) be the resulting EP and BP token
lists, respectively. Let \(V_{\mathrm{ep}}\) and \(V_{\mathrm{bp}}\) denote
their sets of normalized token types, let \(A\) be the selected set of
word-alignment pairs, and let \(V_A\) be the subset of token types that
participate in at least one pair in \(A\). The filter uses the following
coverage score:
\[
s(T_{\mathrm{bp}}, T_{\mathrm{ep}}) =
\frac{\lvert V_A \rvert}
{\lvert V_{\mathrm{bp}} \cup V_{\mathrm{ep}} \rvert},
\]
where the notation follows the definitions above.

The coverage score is then compared with a length-adaptive threshold. If
\(L = \max(\lvert T_{\mathrm{bp}} \rvert, \lvert T_{\mathrm{ep}} \rvert)\),
the threshold is
\[
\tau(L)=
\begin{cases}
0.30, & L \leq 8, \\
0.70, & L \geq 28, \\
0.30 + \dfrac{L-8}{28-8}(0.70-0.30), & \text{otherwise.}
\end{cases}
\]
Short subtitle fragments are allowed lower coverage, while longer pairs must
show stronger bilingual alignment. A separate length-ratio check rejects pairs
where one sentence is implausibly short or long relative to the other. If \(r\) is
the ratio between EP and BP character length, the pair is kept only when
\(0.5 \leq r \leq 2.0\). The final rule therefore keeps pairs that pass both
the alignment-coverage condition and the length-ratio condition.

% Start the complete review procedure on a fresh page instead of leaving a four-line tail.
\clearpage
\section{GPT-Wikipedia}
\label{sec:appendix-gpt-wikipedia}

GPT-Wikipedia is an internally generated resource of near-literal EP/BP pairs
whose sentences are generated through the GPT API using Portuguese Wikipedia
passages as inspiration. After generation, each pair passes through an automatic audit
before human review. The EP and BP sentences are first tokenized, then the
tokens are normalized by lowercasing, removing diacritics, and suppressing
hyphen differences. The two ordered token sequences are compared so that
matching tokens or regions are aligned, while contiguous insertions, deletions,
or replacements become differing spans.

\looseness=-1
The automatic audit then checks these differing spans against a small allowlist
of expected EP/BP correspondences. Remaining differences are flagged for review
and assigned one of three reason categories: unexpected or non-approved lexical
replacements, structural or word-order differences, and reordering artifacts.
Repeated normalized replacements are grouped together and ordered by frequency.
If the same EP expression and BP expression occur in several generated pairs,
the reviewer sees that expression pair once together with its occurrence count
rather than reviewing every sentence pair separately.

Human review labels each grouped replacement as either a valid variety contrast
or an unnecessary change, and high-frequency replacements can be reviewed in
batches. The review scripts can propagate decisions conservatively to related
inflectional forms and containing phrases. When the resulting decision is
applied to matching generated examples, a valid variety contrast preserves the
EP form on the EP side and the BP form on the BP side, whereas an unnecessary
change is equalized across the pair. This applies the recorded and conservatively
propagated decisions consistently without requiring pair-by-pair review of the
complete resource.
